
     \documentclass[fleqn]{article}      
      \usepackage{amsmath}
      \usepackage{fontspec}
      \usepackage{kotex} % 한국어 지원  

      \begin{document}      
       \\\\

쉬움: \\\\
1. 정의 - 행렬은 숫자를 직사각형 형태로 배열한 것입니다. \\\\
2. 예 - 예를 들어, \(\begin{bmatrix} 1 & 2 \\ 3 & 4 \end{bmatrix}\)는 \(2 \times 2\) 행렬입니다. \\\\
3. 덧셈 - 행렬의 덧셈은 같은 크기의 행렬끼리 각 원소를 더하는 것입니다. \\\\
4. 결과 - \(\begin{bmatrix} 6 & 8 \\ 10 & 12 \end{bmatrix}\)입니다. \\\\
5. 의미 - 행렬의 크기는 행렬의 행의 수와 열의 수를 나타냅니다. \\\\

중간: \\\\
1. 조건 - 첫 번째 행렬의 열의 수와 두 번째 행렬의 행의 수가 같아야 합니다. \\\\
2. 결과 - \(\begin{bmatrix} 19 & 22 \\ 43 & 50 \end{bmatrix}\)입니다. \\\\
3. 전치 - \(\begin{bmatrix} 1 & 3 \\ 2 & 4 \end{bmatrix}\)입니다. \\\\
4. 배율 - 스칼라 배율은 행렬의 모든 원소에 같은 숫자를 곱하는 것입니다. \\\\
5. 성질 - 단위 행렬은 곱셈에 대한 항등원입니다. \\\\

어려움: \\\\
1. 역행렬 - 역행렬은 \(\begin{bmatrix} -2 & 1 \\ 1.5 & -0.5 \end{bmatrix}\)입니다. \\\\
2. 조건 - 두 행렬이 모든 원소에서 같을 때입니다. \\\\
3. 행렬식 - 행렬식은 \(-2\)입니다. \\\\
4. 개념 - 고유값은 행렬을 변환할 때 크기만 변화시키는 값입니다. \\\\
5. 설명 - 선형 독립성은 행렬의 열이나 행이 서로 독립적일 때를 말합니다. \\\\      
      \end{document} 
  